\section{Discussion}
\label{sec:discussion}

\subsection{Status of the proposal}

The present work separates \emph{model specification} from \emph{emergence claims}.
The continuum action \eqref{eq:action} defines the substrate dynamics.
The geometric-phase constructions (\Cref{sec:geophase}) define a map from substrate waves to
gauge-like variables. Whether these variables behave like known fields is an empirical and numerical question.

\subsection{Comparison points}

\paragraph{Emergent Lorentz symmetry.}
As in analogue gravity programs, Lorentz-like invariance is expected as an effective symmetry of
the linear-wave cone, not a fundamental symmetry of the microscopic equations \citep{Barcelo2011AnalogueGravity}.
This suggests natural regimes of validity and potential high-frequency deviations.

\paragraph{Geometric phase in classical versus quantum settings.}
Geometric phases are not restricted to quantum mechanics; they arise in classical wave systems and polarization transport
\citep{Pancharatnam1956,Cisowski2022,Bliokh2015SpinOrbit}. This supports the methodological choice of using Berry-type
connections as diagnostics without assuming quantization from the outset.

\paragraph{Ontological models and quantum foundations.}
The substrate hypothesis is closer in spirit to ontological approaches that treat wave structure as physically real
(e.g., de Broglie's early program \citep{deBroglie1924}) while remaining agnostic about whether standard
quantum postulates emerge fully.

\subsection{Limitations and open problems}

Key gaps remain:
\begin{itemize}
\item selecting constitutive laws and parameter regimes that support both stable wave cones and localized modes,
\item deriving clear effective field equations for $A_\mu$ or $F_{\mu\nu}$ (if any) from coarse-graining,
\item demonstrating robust, discretization-independent holonomy behavior,
\item connecting any emergent gauge sector to quantitative electromagnetic phenomenology.
\end{itemize}

The paper therefore prioritizes falsifiable diagnostics and a transparent separation between continuum assumptions
and numerical realizations.
